\begin{lemma}[Situation 7.3: preparation in the disjoint-component case]
Fix a finite loopless multigraph with its individual edge identities, a
palette of \(k\) colors, a distinguished edge \(e_0\), an input coloring,
an exact terminal tree-sequence \(T\), a generating coloring, a
positive-rung ETT, a hierarchy base, a hierarchy counterexample \(C\), and
a \noderef{HierarchyMinimumCounterexampleContext} for these data.  Let
\(P\) be a \noderef{HierarchyReplacementPresentation}, write
\[
 \mu=P.\mathit{coloring},\quad
 H=P.\mathit{counterexample.hierarchy},\quad
 g=P.\mathit{counterexample.good},
\]
and let \(q\) be the last stored index of \(H\).  Put
\[
 j=\noderef{HierarchyPathNumber}(C.\mathit{hierarchy}),\qquad
 \ell=\noderef{HierarchyFullTailLength}(C),
\]
\[
 A=\noderef{HierarchyStarredStage}(H,q),\quad
 B=\noderef{HierarchyFullTailPrefix}(C,j-2),\quad
 X=\noderef{HierarchyFullTailPrefix}(C,j-1),\quad
 Y=\noderef{HierarchyFullTailPrefix}(C,j).
\]
Let \(T_n\) be the terminal ladder tree of \(P\)'s replacement ETT and
define
\[
 D=\noderef{HierarchyStoredColorsAt}(H,q),\qquad
 \Gamma=\bigcup_{\eta\in D}g.\Gamma^q_\eta,
\]
\[
 M_0=\bigcup_{v\in V(P.\mathit{replacementBase.core})\setminus V(T_n)}
       \noderef{MissingColors}(G,\mu,v),
\]
\[
 U_j=\{c:\ \exists r\in\operatorname{Fin}|T|,
 |A|\le r<|A|+j,\ \mu(T.\mathit{edge}(r))=c\}.
\]
Thus \(U_j\) records exact displayed edge identities rather than endpoint
pairs.

Assume \(2\le j<\ell\), and assume that every vertex has a missing color
under every proper coloring leaving exactly \(e_0\) uncolored.  Let
\(\alpha,\beta\) satisfy the actual Claim-\(j-1\) and color-choice
conditions, written inline as follows:
\[
 \alpha\in\noderef{MissingColors}(G,\mu,y_\ell),\qquad
 \alpha\in\noderef{TreeMissingColors}(\mu,B),
\]
and either
\[
 \alpha\notin\Gamma\cup M_0,
\]
or there is \(\eta\in D\) such that
\[
 \alpha\in g.\Gamma^q_\eta,\qquad
 \noderef{HierarchyFirstMissingPosition}(\mu,T,\eta).\mathit{val}
       \le |B|.
\]
Moreover assume
\[
 \beta\ne\alpha,\quad
 \beta\in\noderef{TreeMissingColors}(\mu,B),\quad
 \beta\notin\Gamma,\quad \beta\notin U_j,
\]
and either
\[
 \beta\notin M_0\cup D,
\]
or
\[
 \beta\in D,\qquad
 (g.\Gamma^q_\beta\cup\{\beta\})\cap U_j=\varnothing.
\]
Finally, for
\[
 Q=\noderef{TwoColorComponent}(G,\mu,\alpha,\beta,y_\ell),
\]
assume the literal vertex-disjointness condition
\[
 \operatorname{inl}(v)\notin Q\qquad(v\in V(X)).
\]

Then there exists a
\noderef{HierarchySituationThreeDisjointData} for
\(P,\alpha,\beta\).  In particular, after switching \(Q\), it contains an
actual replacement presentation relative to the fixed original ETT, not
merely one relative to \(P\)'s intermediate replacement ETT.
\end{lemma}
\begin{proof}
Set
\[
 \mathcal E_m=P.\mathit{corresponding.replacementETT},\qquad
 C_m=P.\mathit{counterexample},\qquad
 \mathit{ctx}_m=P.\mathit{minimumContext}.
\]
The \noderef{HierarchyReplacementPresentation} fields put all three
objects over the same replacement input and coloring \(\mu\).  They also
give maximum property, the base of \(H\), and
\noderef{SameGoodHierarchy} from \(C.\mathit{good}\) to \(g\).
Unfolding its \noderef{SameHierarchyStages} component identifies the
divider count, the exact ordered stages, all divider positions, the base
core, and, in the PE case, the base closure.  Its remaining fields
identify every \noderef{HierarchyStoredColorsAt} set and every reserved
Gamma pair.  Consequently the starred stage \(A\), and all four full-tail
prefixes used below, have the same individual vertex and edge entries in
the original and replacement presentations.

Apply \noderef{HierarchyCounterexamplePenultimateElementary} to the six
successive fields of \(\mathit{ctx}_m\): properness,
exact-one-uncoloredness, the full lower-rung package, fewer-divider
induction, rung-local lexicographic minimality, and the \(2n+11\)
estimate.  We obtain
\[
 B_{\rm pen}:=\noderef{HierarchyFullTailPrefix}(C,\ell-1),\qquad
 \noderef{ColoredVertexSetElementary}(G,\mu,V(B_{\rm pen})).       \tag{1}
\]
The definitions of \noderef{HierarchyFullTailLength} and
\noderef{HierarchyFullTailPrefix}, together with the exact nested stages
in \noderef{TashkinovHierarchy} and \(2\le j<\ell\), give
\[
 A\prec B\prec X\prec Y\preccurlyeq B_{\rm pen}\prec T.          \tag{2}
\]
Here each relation is \noderef{TreeSequenceInitialSegment}; in
particular it includes equality of every carried vertex and individual
edge entry.  Natural subtraction causes no endpoint ambiguity because
\(j\ge2\), and \(j<\ell\) puts every displayed length strictly below
\(|T|\).

Let \(z=\operatorname{inl}(y_\ell)\), set
\[
 Q=\noderef{TwoColorComponent}(G,\mu,\alpha,\beta,z),\qquad
 \sigma=\noderef{TwoColorComponentSwap}(G,\mu,\alpha,\beta,z).    \tag{3}
\]
The color-choice assumption gives \(\alpha\ne\beta\).  Since \(\alpha\)
is missing at \(y_\ell\),
\noderef{MissingColorIncidenceEndpoint} says that \(z\) has at most one
neighbor in the \noderef{TwoColorIncidenceGraph}.  The node \(z\) lies in
its own connected component, so it is the endpoint witness in the
definition \noderef{TwoColorPathComponent}.  Thus \(Q\), including the
possible isolated-vertex case, is a certified two-color path component.

We also verify the component-separation premise used by stable change.
Since \(\alpha\in\noderef{TreeMissingColors}(\mu,B)\), expansion of that
definition supplies a displayed \(\alpha\)-missing vertex of \(B\).
The least-candidate definition
\noderef{HierarchyFirstMissingPosition} therefore puts
\[
 v_\alpha=T.\mathit{vertex}
   (\noderef{HierarchyFirstMissingPosition}(\mu,T,\alpha))
       \quad\hbox{in }V(B)\subseteq V(X).                       \tag{4}
\]
The disjointness hypothesis excludes \(\operatorname{inl}(v_\alpha)\)
from \(Q\).  Hence \(Q\) is not the component through that node, and
\noderef{TwoColorComponentsDisjoint} supplies disjointness of the two
full connected components.  This proves the exact separation required
by both orientations of the Cases 1/2 provider.

By \noderef{TwoColorComponentSwapProper}, \(\sigma\) is proper.  The
formula in \noderef{TwoColorComponentSwap} changes only values
\(\mathsf{some}(\alpha)\) and \(\mathsf{some}(\beta)\), and fixes
\(\mathsf{none}\).  Hence the exact uncolored edge set is unchanged;
the \noderef{IsOneEdgeUncolored} field of \(\mathit{ctx}_m\) shows that
exactly \(e_0\) remains uncolored.                                  \(\tag{5}\)

We next prove the two forms of invariance actually needed later.  If
\(v\in V(X)\), then \(\operatorname{inl}(v)\notin Q\).  If an edge-node
of \(Q\) were \noderef{Incident} with \(v\) and colored \(\alpha\) or \(\beta\),
the defining adjacency in \noderef{TwoColorIncidenceGraph} would put
\(\operatorname{inl}(v)\) in the same component, a contradiction.
Thus every edge incident with \(X\) keeps its value in (3).  In
particular every displayed edge identity of \(X\) keeps its value, and
every incident color set, hence every
\noderef{MissingColors} set, is unchanged at a vertex of \(X\).  By the
definition of \noderef{TreeColoringInvariant},
\[
               X\text{ is }(\mu,\sigma)\text{-invariant}.       \tag{6}
\]
The witnesses for \(\alpha,\beta\in
\noderef{TreeMissingColors}(\mu,B)\), transported through (2) and (6),
give
\[
 \alpha,\beta\in\noderef{TreeMissingColors}(\sigma,X).          \tag{7}
\]

A stronger vertexwise equality holds on \(B_{\rm pen}\), even though
the component may meet that larger prefix.  Fix \(v\in V(B_{\rm pen})\).
Colors other than \(\alpha,\beta\) are globally fixed by (3).  If
\(\operatorname{inl}(v)\notin Q\), the adjacency argument above fixes
all incident \(\alpha\)- and \(\beta\)-edges as well.  Suppose instead
\(\operatorname{inl}(v)\in Q\).  The two affected colors already have
missing witnesses in \(B\subseteq B_{\rm pen}\).  By elementarity (1),
neither is missing at \(v\), unless \(v\) is its unique witness vertex;
but every such witness in \(B\) lies in \(X\) and is excluded from \(Q\).
Thus both colors are present at \(v\).  Properness in
\noderef{IsProperEdgeColoring} makes the incident edge of each color
unique.  Their edge-nodes are adjacent to \(\operatorname{inl}(v)\), so
both belong to \(Q\), and (3) exchanges them.  The set of incident colors
is therefore unchanged.  Using the definition of
\noderef{PresentColors} and taking its complement in the palette gives
\[
 \noderef{MissingColors}(G,\sigma,v)
 =\noderef{MissingColors}(G,\mu,v)
          \qquad(v\in V(B_{\rm pen})).                         \tag{8}
\]

We now extract the exact old augmentation routes.  Inspect the terminal
positive certificate of \noderef{TashkinovSeries}.  Its
\noderef{TashkinovSeriesStep} is PE, SE, or RE.  In PE, its
\noderef{TashkinovClosure} supplies a
\noderef{TashkinovAugmentationSequence} from \(T_n\) to the final
certified closure.  In SE or RE it supplies an individual edge \(f_n\),
the literal singleton identity \(F_n=\{f_n\}\), an exact
\noderef{TreeSequenceAdjoinEdge} seed \(T_n+f_n\), and an augmentation
sequence from that seed.  The \(\mathit{belowFinal}\) field of
\noderef{ExtendedTashkinovTree} places \(T\), hence \(X\), in the same
final closure.  The hierarchy and base fields place the relevant PE
source, or exact SE/RE seed, no later than \(A\prec X\).

Every augmentation raises length by one and retains the preceding exact
display by \noderef{TashkinovAugmentation}.  Induction on the finite
reflexive-transitive relation
\noderef{TashkinovAugmentationSequence} therefore locates the unique
members of the closure chain having lengths \(|X|\) and \(|T|\).
Entrywise equality of equal-length initial segments identifies them with
the actual \(X\) and \(T\).  Equivalently, applying
\noderef{TashkinovAugmentationSequenceRestrict} at each cutoff and then
factoring the finite chain yields an old-color suffix
\[
                 X\longrightarrow T\quad\text{under }\mu.      \tag{9}
\]
Now apply
\noderef{TashkinovAugmentationSequenceComponentSwap} to (9), (3), and
the two missing-color witnesses (7).  It retains the same ordered
vertices and individual added edges and gives
\[
                 X\longrightarrow T\quad\text{under }\sigma.  \tag{10}
\]

We record the two universal hierarchy checks needed in every
reconstruction.  Fix \(\eta\in D\).  If
\(\eta\notin\{\alpha,\beta\}\), its missing predicate is globally fixed
by (3), so the finite candidate set and its least element in
\noderef{HierarchyFirstMissingPosition} are unchanged, including the
empty-candidate fallback.  If \(\eta=\alpha\) or \(\eta=\beta\), it has
a missing witness in \(B\), and (8) identifies the entire missing
predicate through \(B_{\rm pen}\).  Both old and new least positions are
therefore at most \(|B|\) and equal; changes after \(B_{\rm pen}\) cannot
alter a least candidate already present.  Hence
\[
 \noderef{HierarchyFirstMissingPosition}(\sigma,T,\eta)
 =\noderef{HierarchyFirstMissingPosition}(\mu,T,\eta)
                       \qquad(\eta\in D).                     \tag{11}
\]
For the Gamma cutoff, \(\beta\notin\Gamma\) makes the implication
vacuous for \(\beta\).  If \(\alpha\notin\Gamma\), it is vacuous for
\(\alpha\) also.  Otherwise the inline Claim-\(j-1\) alternative supplies
an owner \(\eta_0\in D\), membership
\(\alpha\in g.\Gamma^q_{\eta_0}\), and first position at most \(|B|\).
Pairwise disjointness of the Gamma pairs in
\noderef{GoodHierarchy} says that any Gamma pair containing \(\alpha\)
has that same owner.  Therefore, for every \(\eta\in D\),
\[
 (\alpha\in g.\Gamma^q_\eta\text{ or }
   \beta\in g.\Gamma^q_\eta)
 \Longrightarrow
 \noderef{HierarchyFirstMissingPosition}(\mu,T,\eta).\mathit{val}
       \le |B|<|X|.                                           \tag{12}
\]

We next construct a full replacement ETT, base, hierarchy, and goodness
under \(\sigma\), all initially relative to \(\mathcal E_m\).  Define
the literal Cases 1/2 proposition from
\noderef{StableChangeCasesOneTwo} for the ordered pair
\((\alpha,\beta)\), and also for \((\beta,\alpha)\).

Suppose first that one orientation satisfies that proposition.  Apply
\noderef{StableChangeCasesOneTwo} to \(C_m\), the six fields of
\(\mathit{ctx}_m\), the elementarity (1), the two prefix-missing
witnesses, the path certificate above, equation (3), the oriented case
premise, and the component separation following (4).  It gives
\noderef{HierarchyWeaklyStable} and
\noderef{ModColoring} under \(\sigma\).  Swapping the order of the two
colors changes neither \noderef{TwoColorIncidenceGraph}, its component,
nor the formula in \noderef{TwoColorComponentSwap}; these equalities are
obtained by commuting the two displayed disjuncts and the two exchange
branches in those definitions.

Take the split tree to be \(X\).  Equations (2), (6), (7), (8), (9),
(11), and (12) are respectively the strict \(A\)-to-\(X\) prefix, exact
edge fixation, split missingness, penultimate vertex equality, old
suffix, stored-position equality, and affected-Gamma cutoff required by
\noderef{StableChangeCasesOneTwoSplitTail}.  It remains only to give its
literal route from the terminal ladder seed to \(X\).  Return to the
actual terminal \noderef{TashkinovSeriesStep} described before (9).
For positive divider count, weak stability contains
\noderef{TreeColoringInvariant} on \(A\); combine it with (6) and use
\noderef{TashkinovAugmentationSequenceColoringTransport} on the exact
old route through \(X\).  At divider zero in PE,
\noderef{HierarchyBase} supplies the exact
\noderef{OrderedTreeUnion} of \(T_n\) and the PE closure, and
\noderef{CombinedStableClosurePreserved} supplies edge and missing-set
transport on both pieces; (6) transports the remainder through \(X\).
At divider zero in SE/RE the nonparallel base is \(A=T_n\); retain the
individual certified seed \(f_n\), its singleton connecting set, and its
\noderef{TreeSequenceAdjoinEdge} certificate.  Its old endpoint lies in
\(A\subseteq X\), so component disjointness fixes that exact edge, and
(6), (8), and
\noderef{TashkinovAugmentationSequenceColoringTransport} transport the
remaining route.  These are exactly the PE and SE/RE alternatives of the
split-tail theorem, with no identification of parallel edges.

Apply \noderef{StableChangeCasesOneTwoSplitTail}, with source ETT
\(\mathcal E_m\), using those literal premises.  It returns
\[
 \mathcal C_1:
 \noderef{CorrespondingExtendedTashkinovTree}
       (\mathcal E_m,\sigma,T),                                \tag{13}
\]
together with maximum property, a base \(\mathcal B_1\), a hierarchy
\(K_1\), a good witness \(g_1\), and
\noderef{SameGoodHierarchy} from \(g\) to \(g_1\).

For completeness, suppose neither ordered pair satisfies the literal
Cases 1/2 proposition.  We do not add that proposition as an extra
hypothesis, since it is not part of the manuscript Case 2 statement.
Instead use the fact that \(Q\) is wholly outside \(X\).  Apply
\noderef{GoodHierarchyTerminalPrefixRepackage} to (2), obtaining on the
exact tree \(X\) an ETT \(\mathcal E_X\) with the same rung and complete
certified series as \(\mathcal E_m\), maximum property, the same base
core and closure, the retained hierarchy \(H_X\), and a good hierarchy
with the same stored sets and Gamma pairs.  The shared series and the
entrywise equalities construct directly a
\noderef{CorrespondingExtendedTashkinovTree} from
\(\mathcal E_m\) to \(\mathcal E_X\) under \(\mu\).

The one-root set \(\{z\}\), the assumed exclusion of every vertex-node of
\(X\), and equation (3) are precisely an
\noderef{OutsideKempeRecoloring} relative to \(X\).  Apply
\noderef{HierarchyLEMStable} to \(\mathcal E_X,H_X\), using maximum
property, positivity, the lower-rung field of \(\mathit{ctx}_m\), and
the properness and exact-one-uncoloredness fields of that context.  It
gives properness, one-uncoloredness, stability and invariance on \(X\),
and a corresponding maximum ETT with the same good hierarchy under
\(\sigma\).  Restrict its stable-coloring field to the terminal ladder
tree and apply \noderef{Td}; the unchanged certified series identifies
the lower-rung premise and yields
\noderef{ModColoring} of \(\mathcal E_X\) under \(\sigma\).

Independently return to the old certified route to \(X\) extracted before
(9).  Every vertex and edge used before reaching \(X\) lies in \(X\), so
(6), (8), and
\noderef{TashkinovAugmentationSequenceColoringTransport} carry that route
to \(\sigma\), retaining the individual SE/RE seed when present.
Concatenate it with (10) in the reflexive-transitive relation
\noderef{TashkinovAugmentationSequence}.  This gives the exact PE route
\(T_n\to T\), or the exact SE/RE route \(T_n+f_n\to T\), under \(\sigma\).

If the terminal kind is PE, restrict the stability and invariance returned
by \noderef{HierarchyLEMStable} to the common terminal ladder and PE
closure; this is the required \noderef{CombinedStableColoring}.  Apply
\noderef{HierarchySplitter} to \(\mathcal E_X\), its maximum witness,
the lower-rung field, its base and hierarchy, this combined stability,
and the PE route just constructed.  Use the corresponding ETT, maximum
witness, base, and hierarchy returned together by that one application.
If the terminal kind is SE or RE, use the mod-coloring universal clause at
\(T\) with the retained exact seed and route, and apply
\noderef{MaximumPropertyTransport}.  Correspondence identifies the
terminal ladder and kind; the nonparallel field of
\noderef{HierarchyBase} identifies the base core with that ladder, so in
this branch the base and hierarchy can be rebuilt without any PE support
choice.  In either branch, compose the resulting certificate with the
direct certificate \(\mathcal E_m\to\mathcal E_X\) by
\noderef{CorrespondingExtendedTashkinovTreeTrans}.  This gives a
certificate (13) relative to \(\mathcal E_m\), with the correct dependent
input, series, and ETT indices.

Build \(g_1\) on the just-selected branch-coherent base and hierarchy as
follows.  Copy all nonterminal hierarchy stages and divider positions and
take \(T\) as the terminal stage.  The exact nesting follows from (2), and
all nonterminal goodness fields are those returned on \(X\).  For the last
avoidance field, unfold \noderef{HierarchySegmentColors}.  Equation
(11) fixes every stored cutoff.  A new segment color can differ from the
old one only by exchanging \(\alpha\) and \(\beta\) on \(Q\).  The color
\(\beta\) is in no Gamma pair; if \(\alpha\) is in a Gamma pair, (12)
places its owner's cutoff no later than \(X\), before any changed edge.
Thus no reserved color enters the new segment.  The
\noderef{HierarchyStageMissingColors}, unique refresh, parallel base, and
internal minus-closed fields involve only retained stages, base data,
stored sets, and Gamma pairs, so they copy.  This verifies every field of
\noderef{GoodHierarchy} and gives
\noderef{SameGoodHierarchy} from \(g\) to \(g_1\).  Hence both branches
produce the same tuple displayed after (13).

The swap makes \(\beta\) missing at \(y_\ell\).  Indeed, elementarity
(1) and the old \(\beta\)-missing witness in \(B\) show that \(\beta\) is
present at the distinct vertex \(y_\ell\), whereas \(\alpha\) is absent
there.  Properness makes the incident \(\beta\)-edge unique.  Its
edge-node is adjacent to \(z\) and belongs to \(Q\), so (3) changes it to
\(\alpha\); every possible incident affected edge is exchanged in the
same component.  No incident \(\beta\)-edge remains.  Equation (8)
retains a \(\beta\)-missing witness in \(B\).  Since \(B\prec T\), these
are two distinct displayed vertices, and hence \(T\) is not elementary
under \(\sigma\).  Positivity, the maximum witness following (13), and
\(g_1\) now form an actual \noderef{HierarchyCounterexample} \(C_1\).

The certificate (13) is still relative only to the intermediate ETT
\(\mathcal E_m\).  Put \(\mathcal C_0=P.\mathit{corresponding}\) and
apply \noderef{CorrespondingExtendedTashkinovTreeTrans} to
\(\mathcal C_0,\mathcal C_1\).  We obtain
\[
 \mathcal C_{10}:
 \noderef{CorrespondingExtendedTashkinovTree}(\mathit{ett},\sigma,T),
                                                                  \tag{14}
\]
relative to the fixed original ETT, together with equality of its
replacement input and heterogeneous equalities of its dependent series
and ETT with those in (13).  Eliminate these equalities in dependency
order and transport the maximum witness, \(\mathcal B_1,K_1,g_1,C_1\).
This is where the replacement-on-replacement provenance is returned to
the original family.

We construct the transported minimum context explicitly.  Properness
and exact-one-uncoloredness are (5), equivalently the corresponding
fields of (14).  Rung agreement transports the entire lower-rung
six-pack.  The divider and stage equalities in
\noderef{SameGoodHierarchy} transport the fewer-divider assertion.
They also identify the two values of
\noderef{HierarchyCounterexampleMeasure}: both path number and terminal
excess use the same exact stage family.  Thus any smaller same-rung,
equal-divider counterexample for \(C_1\) would contradict the
lexicographic-minimality field of \(\mathit{ctx}_m\), in the sense of
\noderef{HierarchyCounterexampleLess}.  Finally the positive terminal
\noderef{StableColoring} field in correspondence identifies the terminal
ladder missing sets; with rung and ladder-tree equality, it transports
the \(2n+11\) estimate.  These are every field of
\noderef{HierarchyMinimumCounterexampleContext}; denote the result by
\(\mathit{ctx}_1\).

Compose \(P.\mathit{sameGood}:C.\mathit{good}\to g\) first with the
transported \(g\to g_1\) relation.  Transitivity of the exact stage,
stored-set, and Gamma-pair equalities in
\noderef{SameGoodHierarchy} gives \(C.\mathit{good}\to g_1\), in that
order.  The certificate (14), transported maximum property and base,
\(C_1\), this composed same-good witness, and \(\mathit{ctx}_1\) are
exactly an original-relative
\noderef{HierarchyReplacementPresentation}; call it \(P_\sigma\).    \(\tag{15}\)

It remains to transport the color-choice information.  The same-good
and correspondence fields in (15) identify \(A,D,\Gamma\), \(T_n\), and
the base core with their \(P_\sigma\) counterparts.  The base core is an
exact prefix of \(A\preccurlyeq X\), so (6) and
\noderef{TreeSequenceVertexSet} identify all missing sets used in
\(M_0\); hence \(M_0^\sigma=M_0\).  Every exact displayed edge counted by
\(U_j\) before the last one lies in \(X\) and is fixed by (6).  The last
one is the individual edge which adjoins \(y_j\) to an earlier displayed
vertex of \(X\), by \noderef{TreeSequence}.  If its edge-node belonged
to \(Q\), adjacency through that old endpoint would put a vertex-node of
\(X\) in \(Q\), contradicting the hypothesis.  Thus this edge is fixed
as well and
\[
                         U_j^\sigma=U_j.                       \tag{16}
\]
Equations (8), (11), and the same-good Gamma equality now transport the
two alternatives of the input choice literally.  In the first case,
\(\beta\) is outside \(M_0^\sigma,U_j^\sigma,\Gamma_\sigma,D_\sigma\).
In the second, it belongs to \(D_\sigma\), remains missing on \(B\), and
its same Gamma pair together with \(\beta\) is disjoint from the exact
set \(U_j^\sigma\).  This is condition (9), and in both alternatives
\(\beta\notin\Gamma_\sigma\).  Together with the endpoint calculation
and (6)--(10), this supplies every pre-\(\theta\) field of
\noderef{HierarchySituationThreeDisjointData}.

Apply the assumed vertex-missing property to \(P_\sigma.\mathit{coloring}\),
its properness and exactly-one-uncoloredness, and the last vertex \(y_j\)
of \(Y\).  Choose
\[
 \theta\in\noderef{MissingColors}
       (G,P_\sigma.\mathit{coloring},y_j).                     \tag{17}
\]
Apply \noderef{HierarchyCounterexamplePenultimateElementary} once more,
now to all six fields of \(\mathit{ctx}_1\).  It makes
\(B_{\rm pen}\) elementary under the switched coloring.  If
\(\theta\in\Gamma_\sigma\), the avoidance field of
\noderef{GoodHierarchy}, interpreted by
\noderef{HierarchyStageMissingColors}, supplies a
\(\theta\)-missing vertex on the corresponding retained hierarchy stage,
which lies in \(A\subset X\).  By (2), \(y_j\) is a new vertex after
\(X\) and both vertices lie in \(B_{\rm pen}\).  This contradicts
elementarity and (17).  Therefore
\[
                         \theta\notin\Gamma_\sigma.            \tag{18}
\]
The presentation (15), equations (3), (5)--(10), condition (9), and the
chosen color (17)--(18) fill every field of the claimed
\noderef{HierarchySituationThreeDisjointData}.
\end{proof}
